\documentclass[11pt]{article}
\usepackage[T1]{fontenc}
\usepackage[utf8]{inputenc}
\usepackage{lmodern,amsmath,amssymb,microtype}
\usepackage[letterpaper,margin=0.88in]{geometry}
\usepackage[colorlinks=true,urlcolor=blue]{hyperref}
\hypersetup{pdfauthor={Wentao Li},pdftitle={Exactly when variable-overlap permutation packings factor}}
\setlength{\parindent}{0pt}\setlength{\parskip}{4pt}
\newcommand{\done}{\hfill\(\square\)}
\begin{document}
\begin{center}{\Large\bfseries Exactly when variable-overlap\\ permutation packings factor}\par
\small Wentao Li \enspace$\cdot$\enspace Prepared with assistance from OpenAI GPT-6.2 Astra \enspace$\cdot$\enspace 18 November 2026\end{center}

Proposition 6.1 of Blitvi\'c, Kammoun and Steingr\'imsson gives a product formula when occurrences of a permutation overlap in one or two positions. Their Question 2 asks which pattern--overlap combinations admit this generalization. We characterize precisely when a designated overlap splits \emph{every} word of allowed overlaps into independent counting factors. The test returns a short explicit counterexample whenever factorization fails. [\href{https://arxiv.org/html/2610.29846v1\#S6}{Part I, \S6}; \href{https://arxiv.org/html/2610.29846v1\#S7}{\S7, item 2}]

Fix a permutation $\pi$ of length $K$. An admissible overlap $a$, $1\le a<K$, has order-isomorphic prefix and suffix of length $a$. A word $w=a_1\cdots a_r$ specifies $r+1$ copies on consecutive windows with strides $K-a_i$. Let $F(w)$ count permutations on their union with these prescribed occurrences, allowing additional ones. This is the number of linear extensions of the poset generated by the total orders of the individual windows. The empty word gives a single window and $F(\varnothing)=1$.

Fix an alphabet $A$ of admissible overlaps and a seam $m\in A$. Our property is
\[
                   F(u\,m\,v)=F(u)F(v)\qquad(u,v\in A^*).
 \tag{1}
\]
This is equivalent to taking the product over all blocks between successive $m$'s in an arbitrary word. Indeed, apply (1) repeatedly. Conversely, the products for the words with and without a marked split agree block by block.

\subsection*{Interleavings at the shared chain}
The portions on the left and right of the seam intersect in precisely its $m$ positions, which form a chain $s_1<\cdots<s_m$. This holds even if a position occurs in three or more windows: the geometric intersection of the two unions is exactly the overlap of the two adjacent windows. Amalgamating two posets along this common chain preserves each poset's induced order, since any comparison path going out and back across the chain can be shortened between its two chain vertices. It also yields a poset since a cycle would require incompatible orders on the chain.

Two chosen linear extensions on the two sides order their exclusive vertices in each of the $m+1$ gaps around the chain. If they place $a_h,b_h$ vertices in gap $h$, respectively, they admit $\prod_{h=0}^{m}\binom{a_h+b_h}{a_h}$ global extensions. Each interleaving preserves the generating comparisons on both sides and every global extension arises exactly once. In particular the global count is at least the product in (1). It is strictly larger exactly when there is some gap which an extension on \emph{each} side can occupy. Across arbitrary words, the choices of the two sides are independent. It remains to detect all occupiable gaps on a single side.

\subsection*{A finite test on cuts of one window}
Perform the following on each side, starting in its window incident to the seam. Use the ranks of the $m$ suffix entries for the left side, and the $m$ prefix entries for the right. Sort them and add sentinel ranks $0$ and $K+1$. For a gap with consecutive anchors of local ranks $l,u$, record the gap as occupiable immediately if $u>l+1$. Otherwise begin at the cut $c=l$. Explore cuts $c\in\{0,\ldots,K\}$ with the transitions below. Mark a gap occupiable if some walk succeeds.

Choose any $a\in A$. To step outward to a new window, let $\alpha_1<\cdots<\alpha_a$ be the shared ranks in the current window and $\beta_1<\cdots<\beta_a$ those in the new one. For a leftward step these are respectively the sorted prefix and suffix ranks of size $a$. A rightward step exchanges them. Write $\beta_0=0$, $\beta_{a+1}=K+1$. Put $i=|\{j:\alpha_j\le c\}|$. If $\beta_{i+1}>\beta_i+1$, the walk succeeds. Otherwise the new cut is $\beta_i$. Terminate an unsuccessful search when all reachable cuts have been seen.

\textbf{Theorem.} Property (1) holds exactly when the sets of gaps marked occupiable on the two sides are disjoint. On failure there exist counterexample words with $|u|,|v|\le K+1$.

\textit{Proof.} For a chosen gap, a vertex can be placed strictly between its anchors in a linear extension exactly when it is neither at most the lower anchor nor at least the upper one. Indeed, adding the corresponding strict comparisons creates no cycle under those conditions. For the two outer gaps omit the missing anchor's comparison. Thus a successful step really supplies a vertex in an extension of that side, and a cut means that all ranks on one side of it lie at most at the lower anchor and those on the other lie at least at the upper one.

Comparisons between previously attached windows and the new window must pass through their shared chain. At a cut $c$, the greatest rank of the new window forced at most at the lower anchor is therefore $\beta_i$ (zero when none). The least forced at least at the upper anchor is $\beta_{i+1}$ (or $K+1$). All intervening ranks can be put in the gap. This proves the transition and its converse. It also justifies stopping on the first success: an extension of the old part extends over any subsequent window, since extensions agree on their shared chain and can be interleaved there. The finite search therefore detects exactly the achievable gaps. The interleaving formula now gives (1) or its strict failure. A shortest successful walk visits no cut twice before succeeding, so uses at most $K+1$ added windows. Taking successful walks on both sides for a common gap proves the witness bound.\done

\textbf{Example.} For $\pi=125346$, allowed overlaps $\{2,3\}$ and seam $m=2$, the occupiable gaps are $\{0,1\}$ on the left and $\{2\}$ on the right, numbered beginning below the smallest anchor. Thus every sequence factors at its 2-overlaps. A single 3-overlap admits three extensions, so the word $(3,2,3)$ has nine. The test also recovers Proposition 6.1: for $\pi=1\sigma K$ and seam 1 all earlier windows lie at most at the shared letter, and all later windows lie at least at it. The last letters of successive overlapping windows increase, as do their first letters, proving the same conclusion for any admissible alphabet in this case.

\small
\textbf{Verification and status.} The implementation returns actual overlap words witnessing strict inequality. The independent verifier counts linear extensions on the physical positions and checks 8,390 product equalities, 324 counterexamples and the displayed example. Run \texttt{python3 verify.py} from the note's directory. Searches through 18 November 2026, including the \href{https://arxiv.org/abs/2611.12829}{companion paper}, found no earlier necessary-and-sufficient criterion for this variable-overlap product. The note awaits independent review.
\end{document}
